% ============================================================
%  Relatório Final --- CICLOPUS COMBA / SOC Analytics Platform
%  Vítor Hugo Magnus Oliveira --- Semestre 2026/1
% ============================================================
\documentclass[12pt,a4paper]{article}

% ---- Codificação e idioma ----
\usepackage[utf8]{inputenc}
\usepackage[T1]{fontenc}
\usepackage[brazil]{babel}

% ---- Matemática e gráficos ----
\usepackage{amsmath}
\usepackage{graphicx}

% ---- Tabelas de qualidade ----
\usepackage{booktabs}

% ---- Links e metadados ----
\usepackage[colorlinks=true, linkcolor=blue, urlcolor=blue, citecolor=blue]{hyperref}

% ---- Geometria da página ----
\usepackage{geometry}
\geometry{margin=2.5cm}

% ---- Utilitários ----
\usepackage{microtype}   % melhora a hifenização e o espaçamento
\usepackage{parskip}     % parágrafo por espaçamento vertical (sem recuo)
\usepackage{xcolor}      % cores nos avisos

% Evita overfull \hbox em linhas com \texttt{} longos (código/caminhos)
\setlength{\emergencystretch}{3em}

% ============================================================
%  Metadados do documento
% ============================================================
\hypersetup{
  pdftitle  = {Network Attack Detection Analysis},
  pdfauthor = {Vítor Hugo Magnus Oliveira},
}

\title{Network Attack Detection Analysis}
\author{%
  Vítor Hugo Magnus Oliveira (341650) \\
  Thomas Schneider Wiederkehr (342606) \\
  Leonardo Gonzatti Heisler (344125)
}
\date{Ciência de Dados \\ Semestre 2026/1}

% ============================================================
\begin{document}
% ============================================================

\maketitle
\thispagestyle{empty}
\newpage

\tableofcontents
\newpage

% ------------------------------------------------------------
\begin{abstract}
\noindent
Este relatório descreve o projeto \textbf{CICLOPUS COMBA}, uma plataforma de
análise para Centros de Operações de Segurança (SOC) desenvolvida sobre o
dataset CICIDS-2017. O pipeline abrange carregamento e limpeza dos dados,
geração de timestamps sintéticos para simulação operacional, treinamento
comparativo de três classificadores (Decision Tree, Random Forest e XGBoost)
e um dashboard interativo em Streamlit que simula o monitoramento de tráfego
em tempo real. O melhor modelo obtido foi o \textbf{Random Forest},
com \textbf{F1 macro de 0{,}9981} e acurácia de \textbf{99{,}94\%} no
conjunto de teste, demonstrando a viabilidade de um sistema de detecção
automatizada de intrusões baseado em aprendizado de máquina para os cenários
de tráfego benigno, DDoS e PortScan.
\end{abstract}

% ============================================================
\section{Objetivo}
% ============================================================

O projeto tem como objetivo a construção de um \emph{pipeline de ponta a ponta}
para detecção de ataques de rede, simulando as condições operacionais de um
Centro de Operações de Segurança (SOC). Os objetivos específicos são:

\begin{itemize}
  \item Implementar um pipeline reprodutível de ingestão, limpeza e engenharia
        de atributos sobre o dataset CICIDS-2017.
  \item Treinar e comparar três modelos de classificação supervisionada para
        distinguir tráfego benigno de dois cenários de ataque: DDoS e
        PortScan.
  \item Gerar timestamps sintéticos plausíveis que permitam simular análises
        temporais típicas de SOC (timelines, picos de incidente, sazonalidade
        diurna) sem distorcer as proporções reais das classes.
  \item Disponibilizar um dashboard interativo (Streamlit + Plotly) com 6
        páginas funcionais, incluindo um painel operacional de SOC com janela
        temporal deslizante e índice de risco em tempo real.
\end{itemize}

% ============================================================
\section{Dataset}
% ============================================================

\subsection{Origem e composição}

O projeto utiliza o dataset \textbf{CICIDS-2017} (\emph{Canadian Institute
for Cybersecurity Intrusion Detection System 2017}), amplamente empregado
em pesquisas de detecção de intrusão. Foram utilizados três arquivos CSV
da coleção \emph{MachineLearningCVE}:

\begin{itemize}
  \item \texttt{Monday-WorkingHours.pcap\_ISCX.csv}: tráfego exclusivamente
        benigno.
  \item \texttt{Friday-WorkingHours-Afternoon-DDos.pcap\_ISCX.csv}: ataques
        DDoS e tráfego benigno.
  \item \texttt{Friday-WorkingHours-Afternoon-PortScan.pcap\_ISCX.csv}:
        ataques de PortScan e tráfego benigno.
\end{itemize}

Após a concatenação e normalização dos rótulos (que apresentam variações de
grafia entre arquivos), foram mantidas apenas as três classes-alvo:
\texttt{BENIGN}, \texttt{DDoS} e \texttt{PortScan}. O dataset final contém
\textbf{71 atributos} de fluxo de rede (mais o rótulo), extraídos pela
ferramenta CICFlowMeter.

\subsection{Limpeza e pré-processamento}

O módulo \texttt{src/preprocessing.py} aplica as seguintes etapas, na ordem
indicada:

\begin{enumerate}
  \item \textbf{Coerção numérica:} todas as colunas de feature são convertidas
        para numérico; valores não conversíveis tornam-se \texttt{NaN}.
  \item \textbf{Tratamento de infinitos:} valores \texttt{Inf}/\texttt{-Inf}
        (comuns em \emph{Flow Bytes/s}) são substituídos por \texttt{NaN}.
  \item \textbf{Imputação:} \texttt{NaN} restantes preenchidos com a mediana
        da respectiva coluna (robusto a outliers).
  \item \textbf{Remoção de duplicatas} exatas.
  \item \textbf{Remoção de colunas constantes} (variância zero).
  \item \textbf{Codificação de rótulos} via \texttt{LabelEncoder} com ordem
        canônica \texttt{[BENIGN, DDoS, PortScan]}.
\end{enumerate}

\subsection{Geração de timestamps sintéticos}

\textbf{\textcolor{red}{Aviso importante:}} o dataset CICIDS-2017 não possui
uma linha temporal contínua adequada para monitoramento operacional. A coluna
\texttt{Timestamp} adicionada pelo módulo \texttt{src/feature\_engineering.py}
é \textbf{100\% sintética} e existe exclusivamente para simulação de análises
de SOC (timelines, picos, sazonalidade). Ela não representa o tempo real dos
fluxos capturados.

A estratégia de geração de timestamps é parametrizada em \texttt{config.SYNTHETIC\_TIME}
e segue a lógica a seguir:

\begin{itemize}
  \item \textbf{Janela simulada:} 5 dias consecutivos a partir de
        \texttt{2017-07-03} (segunda-feira), com semente aleatória \texttt{42}
        para reprodutibilidade.

  \item \textbf{Tráfego benigno (BENIGN):} distribuído por todos os 5 dias
        segundo um \emph{padrão diurno gaussiano}. A função
        \texttt{\_diurnal\_weights} gera pesos de probabilidade por hora do
        dia (0--23) como uma curva em sino centrada no meio do expediente
        comercial (\texttt{business\_hours = (8, 18)}), com piso mínimo de
        0{,}05 para evitar janelas completamente vazias à noite.

  \item \textbf{Classes de ataque (DDoS / PortScan):} os eventos são
        divididos em duas frações:
        \begin{itemize}
          \item \textbf{85\% concentrados} em janelas de incidente pré-definidas,
                criando picos de atividade plausíveis na timeline.
          \item \textbf{15\% espalhados} como ruído de fundo ao longo de todos
                os dias (maior realismo).
        \end{itemize}

        As janelas de incidente configuradas são:
        \begin{itemize}
          \item DDoS: dia 1 das 10h às 12h; dia 3 das 14h às 16h.
          \item PortScan: dia 0 das 9h às 11h; dia 2 das 13h às 15h;
                dia 4 das 16h às 18h.
        \end{itemize}

  \item \textbf{Colunas derivadas} adicionadas ao DataFrame:
        \texttt{Timestamp}, \texttt{date}, \texttt{hour}, \texttt{weekday},
        \texttt{day\_label}.
\end{itemize}

A feature \texttt{hour} derivada dos timestamps sintéticos é incluída como
atributo de entrada nos modelos (sendo a 71ª feature). Isso implica que
qualquer importância atribuída a \texttt{hour} reflete a distribuição sintética
e não uma característica real dos fluxos (limitação discutida na
Seção~\ref{sec:conclusoes}).

% ============================================================
\section{Metodologia}
% ============================================================

\subsection{Pipeline geral}

O pipeline, orquestrado por \texttt{src/pipeline.py}, executa as seguintes
etapas em sequência:

\begin{enumerate}
  \item Carregamento dos CSVs reais via \texttt{src/data\_loader.py}
        (fallback para dados sintéticos caso os CSVs não estejam presentes).
  \item Amostragem estratificada: caso o dataset ultrapasse
        \texttt{MAX\_SAMPLES = 120.000} registros, aplica-se uma subamostra
        proporcional por classe.
  \item Geração de timestamps sintéticos (\texttt{src/feature\_engineering.py}).
  \item Pré-processamento e codificação (\texttt{src/preprocessing.py}).
  \item Divisão treino/teste estratificada: proporção 75/25,
        \texttt{random\_state = 42}.
  \item Treinamento dos três modelos.
  \item Avaliação e seleção automática do melhor modelo pela métrica
        \textbf{F1 macro} (\texttt{config.SELECTION\_METRIC}).
  \item Serialização do bundle (\texttt{models/models\_bundle.joblib}) e
        exportação das métricas (\texttt{models/metrics.json}).
\end{enumerate}

\subsection{Configuração dos modelos}

Os três classificadores são instanciados com os hiperparâmetros centralizados
em \texttt{config.MODEL\_PARAMS}:

\vspace{0.5em}

\begin{tabular}{@{}lll@{}}
  \toprule
  \textbf{Hiperparâmetro} & \textbf{Valor} & \textbf{Observação} \\
  \midrule
  \multicolumn{3}{@{}l}{\textbf{Decision Tree}} \\
  \quad \texttt{max\_depth} & 18 & Limita overfitting \\
  \quad \texttt{class\_weight} & balanced & Compensa desbalanceamento \\
  \midrule
  \multicolumn{3}{@{}l}{\textbf{Random Forest}} \\
  \quad \texttt{n\_estimators} & 200 & Número de árvores \\
  \quad \texttt{max\_depth} & 22 & Árvores mais profundas que DT \\
  \quad \texttt{class\_weight} & balanced & Compensa desbalanceamento \\
  \quad \texttt{n\_jobs} & $-1$ & Paralelismo total \\
  \midrule
  \multicolumn{3}{@{}l}{\textbf{XGBoost}} \\
  \quad \texttt{n\_estimators} & 300 & Boosting iterations \\
  \quad \texttt{max\_depth} & 8 & Árvores mais rasas (regularização) \\
  \quad \texttt{learning\_rate} & 0{,}2 & Passo de atualização \\
  \quad \texttt{subsample} & 0{,}9 & Subamostragem por iteração \\
  \quad \texttt{colsample\_bytree} & 0{,}9 & Subamostragem de features \\
  \quad \texttt{tree\_method} & hist & Algoritmo aproximado (eficiência) \\
  \quad \texttt{objective} & multi:softprob & Classificação multiclasse \\
  \bottomrule
\end{tabular}

\subsection{Dashboard interativo}

A aplicação \texttt{app/dashboard.py} (Streamlit + Plotly) expõe 6 páginas
navegáveis pela barra lateral:

\begin{enumerate}
  \item \textbf{Overview:} KPIs executivos (total de registros, \% de ataques,
        tráfego benigno, DDoS e PortScan detectados, melhor modelo); gráfico
        de pizza (donut) de proporção por classe e área temporal de volume de
        eventos.
  \item \textbf{Exploração dos Dados (EDA):} distribuição de classes em
        barras; histograma e boxplot interativos por feature (escala log
        opcional); matriz de correlação de Pearson.
  \item \textbf{Monitoramento Temporal:} timeline de eventos por classe com
        granularidade ajustável (30~min/1~h/3~h/6~h); volume por hora do dia;
        ataques por hora; eventos por dia; comparativo DDoS vs.~PortScan;
        tabela de períodos de pico.
  \item \textbf{Feature Importance:} gráfico de barras horizontais com as
        importâncias de atributos do modelo selecionado; tabela completa do
        ranking.
  \item \textbf{Simulador de Predição:} sliders interativos para as 8 features
        mais importantes; presets por perfil de classe (mediana global,
        BENIGN, DDoS, PortScan); exibição da classificação prevista e
        probabilidades por classe.
  \item \textbf{Security Operations Center:} painel com janela temporal
        deslizante (1~h, 2~h, 3~h ou 6~h), navegação por botões e
        \emph{auto-play}; KPIs da janela (eventos, DDoS, PortScan, índice de
        risco 0--100 com nível NORMAL/ELEVADO/CRÍTICO); gauge de risco;
        tabela de alertas ativos; feed dos eventos recentes.
\end{enumerate}

A sidebar disponibiliza filtros globais (classes, período e faixa de horas)
e botão de re-treinamento sob demanda.

% ============================================================
\section{Resultados}
% ============================================================

\subsection{Comparação de desempenho}

Os modelos foram avaliados em um conjunto de teste estratificado composto
por \textbf{30.000 fluxos}: 26.638 BENIGN, 1.500 DDoS e 1.862 PortScan.
A Tabela~\ref{tab:metricas} apresenta as métricas agregadas.

\begin{table}[ht]
  \centering
  \caption{Comparação de métricas entre os classificadores (conjunto de teste).}
  \label{tab:metricas}
  \begin{tabular}{@{}lrrrrr@{}}
    \toprule
    \textbf{Modelo}
      & \textbf{Accuracy}
      & \textbf{Precision}
      & \textbf{Recall}
      & \textbf{F1 macro}
      & \textbf{F1 weighted} \\
    \midrule
    Decision Tree  & 0{,}9991 & 0{,}9970 & 0{,}9978 & 0{,}9974 & 0{,}9991 \\
    \textbf{Random Forest}
                   & \textbf{0{,}9994} & \textbf{0{,}9970} & \textbf{0{,}9991}
                   & \textbf{0{,}9981} & \textbf{0{,}9994} \\
    XGBoost        & 0{,}9993 & 0{,}9974 & 0{,}9987 & 0{,}9980 & 0{,}9993 \\
    \bottomrule
  \end{tabular}
  \small\medskip\\
  \raggedright
  \textit{Precision, Recall e F1 macro = média entre as 3 classes.
  Seleção automática pela métrica F1 macro (\texttt{config.SELECTION\_METRIC}).}
\end{table}

O \textbf{Random Forest} foi selecionado automaticamente como o melhor modelo.
A diferença entre RF (F1 macro 0{,}9981) e XGBoost (0{,}9980) é marginal,
inferior a 0{,}01\% em F1 macro, indicando que ambas as abordagens de
\emph{ensemble} atingem desempenho virtualmente equivalente neste dataset.
A Decision Tree, apesar de ser um classificador mais simples, também alcança
métricas expressivas (F1 macro 0{,}9974), evidenciando a separabilidade
intrínseca das classes no espaço de features de fluxo.

\subsection{Análise por classe: Random Forest}

A Tabela~\ref{tab:por_classe} detalha as métricas individuais do Random Forest
para cada classe do conjunto de teste.

\begin{table}[ht]
  \centering
  \caption{Métricas por classe (Random Forest).}
  \label{tab:por_classe}
  \begin{tabular}{@{}lrrrr@{}}
    \toprule
    \textbf{Classe}
      & \textbf{Precision}
      & \textbf{Recall}
      & \textbf{F1}
      & \textbf{Suporte} \\
    \midrule
    BENIGN    & 0{,}9999 & 0{,}9994 & 0{,}9996 & 26.638 \\
    DDoS      & 0{,}9987 & 0{,}9980 & 0{,}9983 &  1.500 \\
    PortScan  & 0{,}9925 & \textbf{1{,}0000} & 0{,}9963 &  1.862 \\
    \bottomrule
  \end{tabular}
\end{table}

O destaque é o \textbf{Recall = 1{,}0000 para PortScan}: todos os 1.862
casos de varredura de portas foram corretamente identificados, sem nenhum
falso negativo, propriedade crítica para um SOC, pois um PortScan não
detectado pode preceder uma intrusão. A matriz de confusão correspondente é:

\[
\mathbf{C}_{\text{RF}} =
\begin{pmatrix}
  26.622 & 2    & 14 \\
  3      & 1.497 & 0 \\
  0      & 0    & 1.862
\end{pmatrix}
\quad
\begin{array}{l}
  \leftarrow \text{BENIGN (pred.)} \\
  \leftarrow \text{DDoS (pred.)} \\
  \leftarrow \text{PortScan (pred.)}
\end{array}
\]

Os únicos erros observados são 14 fluxos BENIGN classificados como PortScan
(falsos positivos), 2 BENIGN como DDoS e 3 DDoS como BENIGN. São volumes
diminutos num universo de 30.000 amostras.

\subsection{Importância de features (Random Forest)}

A Tabela~\ref{tab:importance} lista as 10 features mais relevantes extraídas
do modelo Random Forest via \texttt{feature\_importances\_}.

\begin{table}[ht]
  \centering
  \caption{Top 10 features por importância (Random Forest, de 71 no total).}
  \label{tab:importance}
  \begin{tabular}{@{}clr@{}}
    \toprule
    \textbf{Rank} & \textbf{Feature} & \textbf{Importância} \\
    \midrule
    1  & Fwd Packet Length Max         & 0{,}0855 \\
    2  & Avg Fwd Segment Size          & 0{,}0814 \\
    3  & Fwd Packet Length Mean        & 0{,}0776 \\
    4  & Subflow Fwd Bytes             & 0{,}0659 \\
    5  & Total Length of Fwd Packets   & 0{,}0579 \\
    6  & Packet Length Mean            & 0{,}0497 \\
    7  & Total Fwd Packets             & 0{,}0361 \\
    8  & Average Packet Size           & 0{,}0337 \\
    9  & Subflow Fwd Packets           & 0{,}0298 \\
    10 & Flow Duration                 & 0{,}0281 \\
    \bottomrule
  \end{tabular}
\end{table}

As features dominantes são relacionadas ao \textbf{tamanho dos pacotes no
sentido direto (forward)} da conexão: \emph{Fwd Packet Length Max},
\emph{Avg Fwd Segment Size} e \emph{Fwd Packet Length Mean} lideram o ranking
com importâncias acima de 0{,}075. Isso é consistente com o comportamento
esperado dos ataques:

\begin{itemize}
  \item \textbf{DDoS:} gera rajadas de pacotes com payloads volumosos
        (\emph{fwd length} alto) e alta cadência.
  \item \textbf{PortScan:} envia pacotes mínimos (SYN isolados), distinguindo-se
        pelo tamanho muito reduzido e pela ausência de tráfego reverso
        (\emph{bwd = 0}).
  \item \textbf{BENIGN:} apresenta fluxos bidirecionais equilibrados com
        payloads de tamanho médio.
\end{itemize}

% ============================================================
\section{Conclusões}
\label{sec:conclusoes}
% ============================================================

\subsection{Interpretação crítica dos resultados}

Os três modelos avaliados atingiram métricas excepcionais (F1 macro $> 0{,}997$),
confirmando que os atributos de fluxo do CICIDS-2017 são altamente discriminativos
para as classes BENIGN, DDoS e PortScan. O \textbf{Random Forest} se destacou
marginalmente, oferecendo o melhor equilíbrio entre Precision, Recall e F1
macro, com a vantagem adicional de Recall perfeito para PortScan,
propriedade de alto valor operacional.

A proximidade entre Random Forest e XGBoost (diferença $< 0{,}01\%$ em F1
macro) é esperada: ambos são modelos de ensemble baseados em árvores, e o
ganho do boosting sobre o bagging é frequentemente superado pela regularização
implícita do bagging em datasets limpos e bem separados.

\subsection{Impacto para um SOC real}

Um sistema com F1 macro próximo a 1{,}000 e Recall = 1{,}0 para PortScan
seria, em princípio, extremamente valioso em produção: detecta virtualmente
todos os ataques com uma taxa de falsos positivos muito baixa. No entanto,
sua adoção em um SOC real exige considerar os pontos a seguir.

\textbf{Pontos favoráveis:}
\begin{itemize}
  \item Inferência em tempo real viável (os modelos de árvore são rápidos e
        serializáveis).
  \item Interpretabilidade via Feature Importance facilita a explicação das
        decisões para analistas de segurança.
  \item O pipeline é reprodutível e modular, permitindo re-treinamento sob
        demanda com novos dados.
\end{itemize}

\textbf{Pontos de atenção:}
\begin{itemize}
  \item \textbf{Generalização para tráfego fora da distribuição:} o CICIDS-2017
        foi capturado em ambiente de laboratório controlado em 2017. Ataques
        modernos podem ter perfis de fluxo distintos dos representados no
        dataset.
  \item \textbf{Dois vetores de ataque apenas:} o modelo atual discrimina
        somente DDoS e PortScan. Um SOC real enfrenta dezenas de categorias
        (exploits, C\&C, exfiltração, ransomware, etc.).
  \item \textbf{Deriva de distribuição (\emph{concept drift}):} o tráfego de
        rede evolui ao longo do tempo; o modelo exigiria re-treinamento periódico
        com dados rotulados recentes.
\end{itemize}

\subsection{Limitações da abordagem de tempo sintético}

A geração de timestamps sintéticos é a limitação mais significativa do projeto
para uso operacional:

\begin{enumerate}
  \item \textbf{Desconexão temporal:} os timestamps não correspondem ao instante
        real de captura dos fluxos. Análises de janela temporal (frequência de
        ataques, sazonalidade) refletem a distribuição escolhida pelo projetista,
        não o comportamento real dos atacantes.

  \item \textbf{Feature \texttt{hour} contaminada:} a 71ª feature (\texttt{hour})
        é derivada do timestamp sintético. Se o modelo aprendeu a associar certas
        horas a certos ataques, essa associação é artificial e pode produzir erros
        caso aplicado a dados com timestamps reais em horários distintos dos
        configurados.

  \item \textbf{Otimismo das métricas:} embora as métricas elevadas sejam em
        grande parte justificadas pela separabilidade real das classes no espaço
        de features, a ausência de variação temporal autêntica pode inflar
        ligeiramente os resultados ao eliminar correlações temporais espúrias
        que ocorrem em tráfego real.
\end{enumerate}

Para um cenário de produção, recomenda-se:
\begin{itemize}
  \item Substituir os timestamps sintéticos pelas marcas de tempo originais dos
        pcaps do CICIDS-2017 (ou de capturas mais recentes).
  \item Excluir a feature \texttt{hour} do conjunto de treino caso os timestamps
        não sejam fidedignos.
  \item Ampliar o escopo para os demais tipos de ataque do CICIDS-2017
        (Brute Force, Web Attacks, Infiltration, Botnet) e adotar um regime de
        validação cruzada temporal (\emph{walk-forward}) ao invés do split
        aleatório.
\end{itemize}

% ============================================================
\section*{Referências}
% ============================================================

\begin{itemize}
  \item Sharafaldin, I., Habibi Lashkari, A., Ghorbani, A.\ A. \textit{Toward
        Generating a New Intrusion Detection Dataset and Intrusion Traffic
        Characterization.} In: ICISSP 2018.
        \url{https://www.unb.ca/cic/datasets/ids-2017.html}

  \item Pedregosa, F.\ et al. \textit{Scikit-learn: Machine Learning in Python.}
        JMLR 12, pp.\ 2825--2830, 2011.

  \item Chen, T., Guestrin, C. \textit{XGBoost: A Scalable Tree Boosting System.}
        KDD 2016.
\end{itemize}

\end{document}
